\section{Robinhood-Style Retail UX}
\label{sec:retail-ux}

The \Zoo{} \DEX{} is engineered to be the simplest, fastest, and
cheapest place for a retail trader to access tokenised global markets.
This section documents the user-facing properties.

\subsection{Zero commission}

Trades on \Zoo{} \DEX{} cost zero commission for retail users. The
fee schedule:

\begin{itemize}[leftmargin=1.2em]
  \item \textbf{Maker rebate.} Resting limit orders that add
        liquidity to the book earn a small rebate funded by the
        taker fee.
  \item \textbf{Taker fee.} Marketable orders pay a small fee
        (default 5 bps for retail, configurable per security class
        by Zoo \DAO{} governance).
  \item \textbf{No commission.} The retail user does not pay a
        per-order commission on top of the spread + taker fee.
        The spread itself is the entry/exit cost.
  \item \textbf{No PFOF (payment for order flow).} The chain is the
        venue; there is no off-chain market-maker paying for the
        right to internalise the order.
\end{itemize}

The taker fee is split between the protocol treasury (for ongoing
protocol development), validator rewards (block-production
incentive), and a conservation allocation (consistent with \Zoo{}
Labs Foundation's 501(c)(3) mission).

\subsection{Fractional shares}

Every listing is fractionable down to one millionth of a share. A
retail user with \$10 to invest can buy fractional positions in any
listed equity, ETF, or REIT. There is no minimum order size other
than the per-trade gas floor, which under \Zoo{} 4.0 \GPU{}-native
gas is sub-cent for typical trade sizes.

Dividend and corporate-action distributions on fractional positions
are computed pro-rata to the holder's exact fractional balance and
distributed automatically by the issuer's action contract.

\subsection{24/7 markets}

The chain runs continuously. Securities trade outside traditional
market hours for instruments whose issuer or transfer agent has
opted into 24/7 trading. For instruments that have not opted in
(e.g., where the underlying market hours are constrained by a
regulator obligation), the chain enforces market hours at the
\Liquidity{} precompile: trades outside the listing's hours window
are rejected.

The matching engine does not pause; only the per-listing eligibility
bit changes. This means a 24/7 instrument continues to match while
a market-hours-constrained instrument is paused, all in the same
block.

\subsection{Instant settlement}

Sub-1.1s end-to-end latency from order submission to triple-cert
finality \cite{lp-020}. Practically:

\begin{itemize}[leftmargin=1.2em]
  \item User taps ``buy''.
  \item Wallet signs and submits to D-Chain sequencer.
  \item Matching engine fills within microseconds.
  \item \Quasar{} 4.0 cert seals within sub-second window.
  \item Wallet shows ``filled, final'' to the user.
\end{itemize}

There is no T+1 / T+2 settlement clearance. The official record at
the transfer agent is updated atomically with the trade. The user's
beneficial ownership of the new position is final at block close.

\subsection{Wallet UX}

The wallet is the user's interface to the \DEX{}. The relevant
panes:

\begin{itemize}[leftmargin=1.2em]
  \item \textbf{Portfolio.} List of positions with current price,
        unrealised P\&L, today's change.
  \item \textbf{Trade.} Symbol search, order entry, order book or
        AMM quote view, order types selector.
  \item \textbf{History.} Filled and unfilled orders with on-chain
        attestation links.
  \item \textbf{Tax.} On-chain cost-basis tracking; per-jurisdiction
        tax-event reports generated from the chain log.
  \item \textbf{Privacy.} Toggle for confidential balances (F-Chain
        FHE) and selective-disclosure proofs (Z-Chain ZKP).
\end{itemize}

The wallet does not run separate market-data infrastructure; it
subscribes to the chain log directly via standard \Lux{} Wallet
RPC (per LP-105 \cite{lp-105}).

\subsection{Mobile parity}

The wallet ships in three form factors --- browser extension,
mobile (iOS/Android), and \Lux{} Desktop --- with the same screen
inventory and identical signing semantics. A trade started on
mobile and finished on desktop signs identically; the chain
treats the order regardless of the form factor of the signer.

\subsection{Make money like the 1\%}

The pitch to retail: \emph{the same instruments and the same
markets that institutional capital trades, on the same venue, with
the same latency, at no commission.} Tokenised treasuries pay the
same yield to a \$50 holder as to a \$50M holder. Fractional REITs
distribute pro-rata. Tokenised perpetuals fund identically per side.
The barrier to access has historically been intermediary fee stacks
and minimum-account-balance rules. \Zoo{} \DEX{}'s zero-commission,
fractional, 24/7, on-chain T+0 settlement removes those barriers
entirely.

\subsection{No dark patterns}

\begin{itemize}[leftmargin=1.2em]
  \item No PFOF: the chain is the venue; nobody pays to internalise
        retail flow.
  \item No spread mark-ups: the spread the user sees is the spread
        the matching engine sees.
  \item No order-routing rebates that misalign with user execution
        quality: the router optimises for user fill price, period.
  \item No gamification of trading frequency: the wallet does not
        push trade volume; defaults are set toward longer-horizon
        positions.
\end{itemize}

This is design-time alignment with consumer protection, not just
post-hoc compliance.

\subsection{Education and risk disclosures}

Every onboarding flow includes risk disclosures appropriate to the
asset class. Derivatives onboarding requires explicit acknowledgement
of leverage and liquidation mechanics. Accreditation-gated
instruments require Z-Chain \ZKP{} proof of accredited status before
the wallet exposes the instrument in search.
